\section{Introducción}

La bioconversión con \textit{Hermetia illucens} transforma residuos
orgánicos en biomasa larvaria y un residuo estabilizado. Desde la
ingeniería ambiental, sin embargo, también es necesario cuantificar el
carbono que retorna a la atmósfera durante el proceso. La producción de
\co{} constituye un observable integrador de la respiración larvaria y
de la degradación microbiana del alimento. Por ello, un modelo útil debe
representar tanto la magnitud total del flujo como su evolución temporal
y su dependencia del sustrato.

El modelo dinámico de Eriksen describe el crecimiento y la respiración
de una larva mediante balances de carbono estructural y lipídico
\cite{eriksenDynamicModellingFeed2022}. Sus reproducciones de los
experimentos de Laganaro y Bekker muestran que el núcleo puede
representar trayectorias individuales de \co{} bajo condiciones
conocidas \cite{laganaroGrowthMetabolicPerformance2021a,
bekkerImpactSubstrateMoisture2021a}. No obstante, esas reproducciones
no constituyen validaciones independientes del presente trabajo, porque
los mismos escenarios intervinieron en la parametrización del modelo de
referencia.

Los experimentos de Ermolaev, Schøn y Rossi aportan condiciones de
transferencia diferentes: medición de emisiones totales, variación de
temperatura y comparación entre tratamientos y controles,
respectivamente \cite{ermolaevGreenhouseGasEmissions2019,
schonEffectTemperatureGrowth2024,
rossiEstimatingDynamicsGreenhouse2024a}. En conjunto, permiten examinar
si un modelo calibrado localmente conserva su capacidad predictiva al
cambiar el alimento, la humedad, la temperatura y el sistema de
observación.

El objetivo de este artículo es formular y evaluar un modelo de caja
gris para la generación de \co{} que conserve el núcleo de Eriksen e
incorpore explícitamente la relación
\[
\text{alimento}\longrightarrow
\text{larvas}+\text{lecho}\longrightarrow \text{\co{}}.
\]
Se distinguen tres preguntas: (i) cuánto del error proviene de aplicar
el núcleo original fuera de su dominio de calibración; (ii) cuánto
aporte representa el lecho microbiano en el total de cámara; y (iii) si
la disponibilidad de alimento permite explicar las diferencias entre
dietas sin sustituir los balances respiratorios originales.

La evaluación separa transferencia sin reajuste, recalibración y
comparación entre configuraciones. Una recalibración con datos propios
mide capacidad de ajuste; no equivale a una validación externa
independiente.
